\chapter{Kirchlich-Zeitgeschichtliches}
\quellseite{32}

\section{I. America}

„The Olive Branch.“ Eine Anzeige dieser Zeitschrift beschließt der „Missionary“ von Pittsburg in der Nummer vom 28. Januar mit der vortrefflichen Bemerkung: „Möge er (der Olive Br.) sich als einen wahren Oelzweig der Kirche erweisen, nicht durch Ignorirung der großen Fragen, welchen er begegnen und die er entscheiden muß, sondern durch ernstliche und liebevolle Arbeit für die Einheit derselben im Glauben und im Leben, und so durch Sicherstellung des Friedens, der Gott allein wohlgefällig ist.“

Der „Lutheran Observer“ vom 29. Januar enthält Worte des Abschiedes von Seiten des bisherigen Redakteurs, des Herrn Dr. Benj. Kurtz, von seinen Lesern.

Eine neue Zeitschrift. Aus dem „Histor. Zeitblatt“ ersehen wir, daß vor kurzem die erste Nummer eines „Kirchenblattes der ev.-luth. Synode von Iowa“, redigirt von den Pastoren G. Großmann und G. Fritschel, erschienen ist.

Der deutsche Kirchenfreund. In dem Decemberheft des vorigen Jahrgangs dieser Zeitschrift findet sich die Fortsetzung einer Biographie Calvin’s. Darin heißt es u.\,A.: „Schon in der ersten Ausgabe seiner Institution … spricht er (Calvin) zwar von einem wirklichen Genuß des Leibes und Blutes Christi im h. A., aber wie er sich gegen die bloß symbolische Ansicht Zwingli’s erklärt, so, und noch viel stärker, spricht er sich gegen die räumliche Gegenwart Luther’s aus. Melanchthon … hatte … das vermittelnde System des Synergismus aufgestellt.“ Darauf folgen denn die alten abgedroschenen Mährlein, nach welchen Luther, nachdem er bereits Calvin’s abweichende Lehre erkannt hatte, denselben dennoch für den Mann angesehen und erklärt haben soll, der „dem Papstthum vollends den Stoß geben und was er (Luther) angefangen, mit Gottes Hülfe hinausführen werde“, ja, dem er (Luther) „hätte anfänglich wohl dürfen die ganze Sache von diesem Streit (über das h. A.) anheimstellen“ — während es documentarisch zu erweisen und von uns bereits vor 12 Jahren erwiesen worden ist und von dem Schreiber der bezeichneten Biographie zugestanden wird, daß Melanchthon nicht einmal „das Herz gehabt“ hat, den einzigen von Calvin an Luther je abgesendeten Brief letzterem zu überreichen! — Über jene Biographie sagt aber nichts desto weniger der „lutherische“ Herausgeber des „Kirchenfreundes“ Folgendes: „Namentlich danke ich hiermit in meinem und Anderer Namen dem Bruder, der uns das Bild des Genfer Reformators in klaren strengen Zügen gezeichnet hat, für seine eingehende Arbeit, die wir zu den Stücken aus den Jahrgängen des Kirchenfreundes zählen, denen um ihres bedeutenden Gegenstandes und instructiven Charakters willen ein bleibender Werth zukommt.“ Wir meinen, ein solches Urtheil bedarf keines Commentars.

\section{II. Ausland}

Copulation. Aus Thüringen bringen die H. M. zur weiteren Ergänzung der Eisenacher Kirchenconferenzbeschlüsse diejenigen, welche in Bezug auf Trauungsverweigerungen von Geistlichen gefaßt sind und darauf hinauskommen, daß mit den Gewissensbedenken der sich weigernden Geistlichen milde verfahren, ihnen aber zur Pflicht gemacht werden solle, nicht nach eigenem Gutdünken, sondern nach den Entscheidungen der kirchlichen Behörden zu verfahren, welche die Prüfung und Entscheidung der Ehegeschiedenen in ihre Hand nehmen würden; also wesentlich, wie in Preußen bereits angeordnet ist. Nur Weimar, Gotha, Altenburg, Sondershausen, Waldeck, Reuß und Hessen-Homburg hatten die unbedingte Verpflichtung des Geistlichen, bürgerlich rechtgültig Geschiedene zu trauen, aufrecht erhalten wollen.
